% GENERATED FILE. Edit canonical lessons, then run npm run generate:books.
% canonical-source-hash: fnv1a64:c86c0eb65c93ddd5
% canonical-lessons: BN-C380-onukul, BN-C380-protikul, BN-C380-upojukto, BN-C380-jogyo, BN-C380-obhyosto-hooya

\chapter{The Right Conditions}
\label{ch:bn-the-right-conditions}
\hlchaptermodality{380}

\begin{chapteropening}
\textbf{By the end of this chapter:} \emph{I can say when conditions are right and what to do if they are not.}

Complete the last lesson of chapter 380: I can say when conditions are right and what to do if they are not.
\end{chapteropening}

\section[\texorpdfstring{\bn{অনুকূল} (ônukūl)}{অনুকূল (ônukūl)}]{\bn{অনুকূল} (ônukūl) — favourable}

\label{lesson:BN-C380-onukul}



\begin{quote}
Before the new one: say the Bengali for to break a promise, then the Bengali for to vow.
\end{quote}

\subsection*{You'll want to know: \bn{অনুকূল}}
\textbf{\bn{অনুকূল}} — \emph{ônukūl} — \textquotedblleft{}favourable\textquotedblright{}.

\begin{quote}
\bn{আবহাওয়া} \bn{অনুকূল} \bn{হলে} \bn{আমরা} \bn{কাল} \bn{নৌকায়} \bn{যাব।}
\end{quote}

\begin{quote}
\emph{ābôhāoyā ônukūl hôle āmrā kāl noukāy jābô.}
\end{quote}

\begin{quote}
If the weather is favourable, we'll go by boat tomorrow.
\end{quote}

\begin{cousinweb}
\textbf{\bn{অনুকূল}} (\emph{ônukūl}) is Sanskrit \emph{anukūla}, 'along the bank', and so going with you, favourable (Monier-Williams).
\end{cousinweb}

\begin{grammarlens}[title={What you've built}]
One more word for saying what happens if.
\end{grammarlens}

\subsection*{Your turn}
\begin{itemize}
\raggedright
  \item \emph{Say it:} \emph{ônukūl}
  \item \emph{Say it:} \emph{ônukūl}, once more
  \item \emph{Say it:} \emph{ābôhāoyā ônukūl hôle āmrā kāl noukāy jābô.} — If the weather is favourable, we'll go by boat tomorrow.
  \item \emph{Recall:} say the Bengali for to break a promise, then the Bengali for to vow, then say \emph{ônukūl} again
\end{itemize}

\begin{tcolorbox}[breakable,colback=teal!4,colframe=teal!35!black,title={Before you move on}]
What does \bn{অনুকূল} mean? (\textquotedblleft{}Favourable\textquotedblright{}.) Say it once more.
\end{tcolorbox}
\section[\texorpdfstring{\bn{প্রতিকূল} (prôtikūl)}{প্রতিকূল (prôtikūl)}]{\bn{প্রতিকূল} (prôtikūl) — unfavourable; adverse}

\label{lesson:BN-C380-protikul}



\begin{quote}
Before the new one: say the Bengali for to vow, then the Bengali for favourable.
\end{quote}

\subsection*{You'll want to know: \bn{প্রতিকূল}}
\textbf{\bn{প্রতিকূল}} — \emph{prôtikūl} — \textquotedblleft{}unfavourable; adverse\textquotedblright{}.

\begin{quote}
\bn{পরিস্থিতি} \bn{প্রতিকূল} \bn{হলে} \bn{আমরা} \bn{ফিরে} \bn{আসব।}
\end{quote}

\begin{quote}
\emph{pôristhiti prôtikūl hôle āmrā phire āsbô.}
\end{quote}

\begin{quote}
If the situation turns against us, we'll come back.
\end{quote}

\begin{cousinweb}
\textbf{\bn{প্রতিকূল}} (\emph{prôtikūl}) is Sanskrit \emph{pratikūla}, 'against the bank', contrary (Monier-Williams), the opposite of \bn{অনুকূল}.
\end{cousinweb}

\begin{grammarlens}[title={What you've built}]
One more word for saying what happens if.
\end{grammarlens}

\subsection*{Your turn}
\begin{itemize}
\raggedright
  \item \emph{Say it:} \emph{prôtikūl}
  \item \emph{Say it:} \emph{prôtikūl}, once more
  \item \emph{Say it:} \emph{pôristhiti prôtikūl hôle āmrā phire āsbô.} — If the situation turns against us, we'll come back.
  \item \emph{Recall:} say the Bengali for to vow, then the Bengali for favourable, then say \emph{prôtikūl} again
\end{itemize}

\begin{tcolorbox}[breakable,colback=teal!4,colframe=teal!35!black,title={Before you move on}]
What does \bn{প্রতিকূল} mean? (\textquotedblleft{}Unfavourable; adverse\textquotedblright{}.) Say it once more.
\end{tcolorbox}
\section[\texorpdfstring{\bn{উপযুক্ত} (upôjuktô)}{উপযুক্ত (upôjuktô)}]{\bn{উপযুক্ত} (upôjuktô) — suitable; fitting}

\label{lesson:BN-C380-upojukto}



\begin{quote}
Before the new one: say the Bengali for favourable, then the Bengali for adverse.
\end{quote}

\subsection*{You'll want to know: \bn{উপযুক্ত}}
\textbf{\bn{উপযুক্ত}} — \emph{upôjuktô} — \textquotedblleft{}suitable; fitting\textquotedblright{}.

\begin{quote}
\bn{উপযুক্ত} \bn{লোক} \bn{পেলে} \bn{আমরা} \bn{কাজটা} \bn{দেব।}
\end{quote}

\begin{quote}
\emph{upôjuktô lok pele āmrā kājṭā debô.}
\end{quote}

\begin{quote}
If we find a suitable person, we'll give them the job.
\end{quote}

\begin{cousinweb}
\textbf{\bn{উপযুক্ত}} (\emph{upôjuktô}) is Sanskrit \emph{upayukta}, suitable or appropriate (Monier-Williams).
\end{cousinweb}

\begin{grammarlens}[title={What you've built}]
One more word for saying what happens if.
\end{grammarlens}

\subsection*{Your turn}
\begin{itemize}
\raggedright
  \item \emph{Say it:} \emph{upôjuktô}
  \item \emph{Say it:} \emph{upôjuktô}, once more
  \item \emph{Say it:} \emph{upôjuktô lok pele āmrā kājṭā debô.} — If we find a suitable person, we'll give them the job.
  \item \emph{Recall:} say the Bengali for favourable, then the Bengali for adverse, then say \emph{upôjuktô} again
\end{itemize}

\begin{tcolorbox}[breakable,colback=teal!4,colframe=teal!35!black,title={Before you move on}]
What does \bn{উপযুক্ত} mean? (\textquotedblleft{}Suitable; fitting\textquotedblright{}.) Say it once more.
\end{tcolorbox}
\section[\texorpdfstring{\bn{যোগ্য} (jogyô)}{যোগ্য (jogyô)}]{\bn{যোগ্য} (jogyô) — qualified; deserving}

\label{lesson:BN-C380-jogyo}



\begin{quote}
Before the new one: say the Bengali for adverse, then the Bengali for suitable.
\end{quote}

\subsection*{You'll want to know: \bn{যোগ্য}}
\textbf{\bn{যোগ্য}} — \emph{jogyô} — \textquotedblleft{}qualified; deserving\textquotedblright{}.

\begin{quote}
\bn{তুমি} \bn{যদি} \bn{যোগ্য} \bn{হও}, \bn{তাহলে} \bn{চাকরিটা} \bn{পাবে।}
\end{quote}

\begin{quote}
\emph{tumi jôdi jogyô hôo, tāhôle chākriṭā pābe.}
\end{quote}

\begin{quote}
If you're qualified, you'll get the job.
\end{quote}

\begin{cousinweb}
\textbf{\bn{যোগ্য}} (\emph{jogyô}) is Sanskrit \emph{yogya}, fit for the yoke and so fit for a task (Monier-Williams). A person is \bn{যোগ্য} for a post or a prize.
\end{cousinweb}

\begin{grammarlens}[title={What you've built}]
One more word for saying what happens if.
\end{grammarlens}

\subsection*{Your turn}
\begin{itemize}
\raggedright
  \item \emph{Say it:} \emph{jogyô}
  \item \emph{Say it:} \emph{jogyô}, once more
  \item \emph{Say it:} \emph{tumi jôdi jogyô hôo, tāhôle chākriṭā pābe.} — If you're qualified, you'll get the job.
  \item \emph{Recall:} say the Bengali for adverse, then the Bengali for suitable, then say \emph{jogyô} again
\end{itemize}

\begin{tcolorbox}[breakable,colback=teal!4,colframe=teal!35!black,title={Before you move on}]
What does \bn{যোগ্য} mean? (\textquotedblleft{}Qualified; deserving\textquotedblright{}.) Say it once more.
\end{tcolorbox}
\section[\texorpdfstring{\bn{অভ্যস্ত} \bn{হওয়া} (ôbhyôstô hôoyā)}{অভ্যস্ত হওয়া (ôbhyôstô hôoyā)}]{\bn{অভ্যস্ত} \bn{হওয়া} (ôbhyôstô hôoyā) — to get used to}

\label{lesson:BN-C380-obhyosto-hooya}



\begin{quote}
Before the new one: say the Bengali for suitable, then the Bengali for qualified.
\end{quote}

\subsection*{You'll want to know: \bn{অভ্যস্ত} \bn{হওয়া}}
\textbf{\bn{অভ্যস্ত} \bn{হওয়া}} — \emph{ôbhyôstô hôoyā} — \textquotedblleft{}to get used to\textquotedblright{}.

\begin{quote}
\bn{এখানে} \bn{এক} \bn{মাস} \bn{থাকলে} \bn{তুমি} \bn{গরমে} \bn{অভ্যস্ত} \bn{হয়ে} \bn{যাবে।}
\end{quote}

\begin{quote}
\emph{ekhāne ek mās thākle tumi gôrôme ôbhyôstô hôye jābe.}
\end{quote}

\begin{quote}
If you stay here a month, you'll get used to the heat.
\end{quote}

Past (he/she …): \bn{তিনি} \bn{নতুন} \bn{শহরে} \bn{অভ্যস্ত} \bn{হলেন}, \emph{tini nôtun shôhôre ôbhyôstô hôlen}.

\begin{cousinweb}
\textbf{\bn{অভ্যস্ত} \bn{হওয়া}} (\emph{ôbhyôstô hôoyā}) is Sanskrit \emph{abhyasta}, practised (Monier-Williams), with \bn{হওয়া}: to become practised at living with something, to get used to it. What you get used to takes -\bn{এ}.
\end{cousinweb}

\begin{grammarlens}[title={What you've built}]
One more word for saying what happens if.
\end{grammarlens}

\subsection*{Your turn}
\begin{itemize}
\raggedright
  \item \emph{Say it:} \emph{ôbhyôstô hôoyā}
  \item \emph{Say it:} \emph{ôbhyôstô hôoyā}, once more
  \item \emph{Say it:} \emph{ekhāne ek mās thākle tumi gôrôme ôbhyôstô hôye jābe.} — If you stay here a month, you'll get used to the heat.
  \item \emph{Recall:} say the Bengali for suitable, then the Bengali for qualified, then say \emph{ôbhyôstô hôoyā} again
\end{itemize}

\begin{tcolorbox}[breakable,colback=teal!4,colframe=teal!35!black,title={Before you move on}]
What does \bn{অভ্যস্ত} \bn{হওয়া} mean? (\textquotedblleft{}To get used to\textquotedblright{}.) Say it once more.
\end{tcolorbox}
